\emph{Symbols.} The statements of (i), (ii) and the first-order part of (iv) are identities between constant-coefficient differential expressions, linear in $h$ or bilinear in $h_{1}, h_{2}$. Such an identity holds for all fields exactly when its symbol identity holds for all vectors: put $p_{1}$ for each derivative of $h_{1}$ and $p_{2}$ for each derivative of $h_{2}$ and regard $h_{1}, h_{2}$ as constant symmetric maps. The symbol of $\mathcal{R}[h]$ is the stored $\mathcal{R}_{p}[h](B) = -\tfrac12\,p \wedge h(B \cdot p)$ and that of $\mathcal{E}[h]$ the stored Einstein map $\mathcal{G}_{p}$ of Theorem \claimref{synthesis.first-order-gravitational-field-has-two-free-directions-at-null-momentum-and-four-constraints}, by the definition of both on the field $\tfrac12(p \cdot x)^{2}h_{0}$ in Theorem \claimref{synthesis.weyl-flat-variations-carry-only-the-stress-trace-and-the-whole-stress-clause-is-the-trace-free-channel}. Contracting $\mathcal{R}_{p}$: with $(e_{b} \wedge a) \cdot p = e_{b}(a \cdot p) - a\,(e_{b} \cdot p)$ and $e^{b} \cdot (p \wedge v) = (e^{b} \cdot p)v - (e^{b} \cdot v)p$, the Ricci map is $\mathrm{Ric}_{p}[h](a) = \tfrac12\bigl(Q(p)\,h(a) + (a \cdot p)\,\tau\,p - (a \cdot p)\,h(p) - (a \cdot h(p))\,p\bigr)$, with trace $Q(p)\tau - p \cdot h(p)$, so
\[
\mathcal{E}[h](a) = \tfrac12\bigl(\nabla^{2}h(a) + (a \cdot \nabla)\nabla\tau - (a \cdot \nabla)d - \nabla(a \cdot d)\bigr) - \tfrac12\,a\,\bigl(\nabla^{2}\tau - \nabla \cdot d\bigr).
\]

\emph{(i).} Antisymmetry is by inspection: each of the five groups changes sign under exchange of the two fields. For the divergence take the groups of $\ip{\mathcal{E}[h_{1}]}{h_{2}} - \ip{h_{1}}{\mathcal{E}[h_{2}]}$ in turn. The terms with $\nabla^{2}h$ give $\tfrac12\nabla \cdot \bigl(\dot\nabla\ip{\dot h_{1}}{h_{2}} - \dot\nabla\ip{h_{1}}{\dot h_{2}}\bigr)$. Write $H_{\tau}$ for the map $a \mapsto (a \cdot \nabla)\nabla\tau$. The terms with $H_{\tau}$, together with those with $a\,\nabla \cdot d$, give $\tfrac12\bigl(\ip{H_{\tau_{1}}}{h_{2}} - \tau_{1}\nabla \cdot d_{2}\bigr) + \tfrac12\bigl(\tau_{2}\nabla \cdot d_{1} - \ip{h_{1}}{H_{\tau_{2}}}\bigr)$, which is $\tfrac12\nabla \cdot \bigl(h_{2}(\nabla\tau_{1}) - \tau_{1}d_{2}\bigr) + \tfrac12\nabla \cdot \bigl(\tau_{2}d_{1} - h_{1}(\nabla\tau_{2})\bigr)$, the cross terms $\nabla\tau_{1} \cdot d_{2}$ and $d_{1} \cdot \nabla\tau_{2}$ cancelling within each bracket. For the terms with $(a \cdot \nabla)d + \nabla(a \cdot d)$: by the symmetry of $h_{2}$ the two halves pair equally with it, and the pairing of $-\tfrac12$ times that map with $h_{2}$ is $-\dot d_{1} \cdot h_{2}(\dot\nabla) = -\nabla \cdot \bigl(h_{2}(d_{1})\bigr) + d_{1} \cdot d_{2}$, and likewise with the fields exchanged; the difference is $\nabla \cdot \bigl(h_{1}(d_{2}) - h_{2}(d_{1})\bigr)$, the terms $d_{1} \cdot d_{2}$ cancelling. The terms with $a\,\nabla^{2}\tau$ give $-\tfrac12(\tau_{2}\nabla^{2}\tau_{1} - \tau_{1}\nabla^{2}\tau_{2}) = -\tfrac12\nabla \cdot (\tau_{2}\nabla\tau_{1} - \tau_{1}\nabla\tau_{2})$. The sum is $\nabla \cdot \omega$. The check verifies the symbol identity $(p_{1} + p_{2}) \cdot \omega = \ip{\mathcal{G}_{p_{1}}[h_{1}]}{h_{2}} - \ip{h_{1}}{\mathcal{G}_{p_{2}}[h_{2}]}$ on twelve pairs of generic symmetric maps and momenta, with both sides nonzero, using the stored curvature formula and not the formula for $\mathcal{E}$ above; at equal momenta it is the symmetry of $\mathcal{G}_{p}$ for the pairing, clause (ii) of the cited theorem.

\emph{(ii).} A displacement changes $h$ at momentum $p$ by $a \mapsto (a \cdot p)\xi + (a \cdot \xi)p$ (the gauge space $G_{p}$ of the cited theorems), so $h(n)$ by $(n \cdot p)\xi + (n \cdot \xi)p$. If $n \cdot p \neq 0$ and this vanishes, contracting with $n$ gives $2(n \cdot p)(n \cdot \xi) = 0$, so $n \cdot \xi = 0$ and then $\xi = 0$: the map is injective, of rank four, and one displacement removes $h(n)$. If $n \cdot p = 0$ the map is $\xi \mapsto (n \cdot \xi)p$, of rank one. Now let $h_{1}(n) = h_{2}(n) = 0$ and contract $\omega$ with $n$. Since the maps are symmetric, $n \cdot h_{2}(\nabla\tau_{1}) = h_{2}(n) \cdot \nabla\tau_{1} = 0$, $n \cdot h_{1}(d_{2}) = 0$ and likewise with the fields exchanged; and $n \cdot d = n \cdot \dot h(\dot\nabla) = \dot h(n) \cdot \dot\nabla = 0$. What is left is $n \cdot \omega = \tfrac12(\ip{\partial_{n}h_{1}}{h_{2}} - \ip{h_{1}}{\partial_{n}h_{2}}) - \tfrac12(\tau_{2}\partial_{n}\tau_{1} - \tau_{1}\partial_{n}\tau_{2})$. A symmetric map with $h(n) = 0$ is $h = h^{TT} + \tfrac12\tau\,\theta_{\perp} + (a \mapsto (a \cdot n)w + (a \cdot w)n)$, with $\theta_{\perp}(a) = -(a \cdot e_{1})e_{1} - (a \cdot e_{2})e_{2}$ the projection on the transverse plane ($\tr\theta_{\perp} = \ip{\theta_{\perp}}{\theta_{\perp}} = 2$) and $w$ in the span of $n, e_{1}, e_{2}$; the last part pairs to zero with every map annihilating $n$, since $\ip{a \mapsto (a \cdot n)w + (a \cdot w)n}{h'} = 2\,w \cdot h'(n)$, and has trace $2\,n \cdot w = 0$. The three maps $\epsilon_{+}, \epsilon_{\times}, \theta_{\perp}$ are orthogonal for the pairing, so $\ip{h_{1}}{h_{2}} = \ip{h_{1}^{TT}}{h_{2}^{TT}} + \tfrac12\tau_{1}\tau_{2}$ and the trace terms combine to $\tfrac14 - \tfrac12 = -\tfrac14$: the stated formula. For the equation along the generators, $Q(n) = 0$ gives $n \cdot \mathcal{G}_{p}[h](n) = n \cdot \mathrm{Ric}_{p}[h](n) = \tfrac12\bigl(Q(p)\,n \cdot h(n) + (n \cdot p)^{2}\tau - 2(n \cdot p)\,n \cdot h(p)\bigr) = \tfrac12(n \cdot p)^{2}\tau$ in the gauge. On solutions $\partial_{n}(\tau_{2}\partial_{n}\tau_{1} - \tau_{1}\partial_{n}\tau_{2}) = \tau_{2}\partial_{n}^{2}\tau_{1} - \tau_{1}\partial_{n}^{2}\tau_{2} = 0$, and an affine function periodic along a generator is constant, with zero derivative. The remaining form $\tfrac12\int(\ip{\partial_{n}h_{1}^{TT}}{h_{2}^{TT}} - \ip{h_{1}^{TT}}{\partial_{n}h_{2}^{TT}})$ contains no derivative across generators, so it is a sum over generators of forms in the data of each; on the two components the pairing is $2(h_{+}h_{+}' + h_{\times}h_{\times}')$, positive. That the flux through $N$ is the whole symplectic form of the free field is the light-front setting, taken here as for the stored photon and not proved. The check verifies the flux formula, in symbols, on twelve pairs of generic maps of the gauge, that adding parts along $n$ changes nothing, the two ranks, and the equation along the generators.

\emph{(iii).} Everything is from clauses (b), (c) and the torsion equation of the Definition; its field equation is not used: $D_{a}M = a \cdot \nabla M + \Omega(a) \times M$ and $A \cdot D = D_{\hb(A)}$ (clause (b)), $[D_{a}, D_{b}]M = R(a \wedge b) \times M$ for constant $a, b$ with $\mathcal{R}(B) = -R(\hb(B))$ (clause (c)), and the torsion equation (clause (f)). Since $\Omega(a) \times$ is a derivation of the product and $(W \times X) \cdot Y + X \cdot (W \times Y) = 0$ for a bivector $W$ and vectors $X, Y$, the derivative respects the product of vectors: $a \cdot \nabla(X \cdot Y) = (D_{a}X) \cdot Y + X \cdot (D_{a}Y)$. For a bivector $\mathcal{B}$ and vectors $x, y$: $x \cdot (\mathcal{B} \cdot y) = -y \cdot (\mathcal{B} \cdot x) = (y \wedge x) \cdot \mathcal{B}$.

\emph{Three lemmas.} (1) For physical vectors $A, C$ at a point and a vector field $X$, $(A \wedge C) \cdot (D \wedge X) = A \cdot (C \cdot D\,X) - C \cdot (A \cdot D\,X)$: indeed $D \wedge X = \partial_{E} \wedge (E \cdot D\,X)$, and $(A \wedge C) \cdot (u \wedge v) = (A \cdot v)(C \cdot u) - (A \cdot u)(C \cdot v)$. (2) For physical vector fields $A, C$ with background images $a = \hb(A)$, $c = \hb(C)$, $\hb(A \cdot D\,C - C \cdot D\,A) = (a \cdot \nabla)c - (c \cdot \nabla)a$. Proof: the torsion equation at the constant vector $e^{\mu}$ gives $D \wedge g^{\mu} = 0$ for $g^{\mu} = \hb^{T}(e^{\mu})$, so by (1) $A \cdot (C \cdot D\,g^{\mu}) = C \cdot (A \cdot D\,g^{\mu})$. The scalar $C \cdot g^{\mu} = c \cdot e^{\mu}$ has derivative $(a \cdot \nabla)(c \cdot e^{\mu}) = (A \cdot D\,C) \cdot g^{\mu} + C \cdot (A \cdot D\,g^{\mu})$; subtracting the same with $A$ and $C$ exchanged, $\bigl((a \cdot \nabla)c - (c \cdot \nabla)a\bigr) \cdot e^{\mu} = (A \cdot D\,C - C \cdot D\,A) \cdot g^{\mu} = \hb(A \cdot D\,C - C \cdot D\,A) \cdot e^{\mu}$ for each $\mu$. (3) For physical vector fields $A, C$ and a multivector field $X$, $[A \cdot D, C \cdot D]X = -\mathcal{R}(A \wedge C) \times X + (A \cdot D\,C - C \cdot D\,A) \cdot D\,X$. Proof: $A \cdot D(C \cdot D\,X) = D_{a}(D_{c}X)$ with $c$ a field; $D_{c}$ is linear in $c$, so this is $D_{a}D_{c}X$ with $c$ held fixed plus $D_{(a \cdot \nabla)c}X$. Hence $[A \cdot D, C \cdot D]X = R(a \wedge c) \times X + D_{(a \cdot \nabla)c - (c \cdot \nabla)a}X$, and by (2) the last term is the derivative along the physical vector $A \cdot D\,C - C \cdot D\,A$, while $R(a \wedge c) = R(\hb(A \wedge C)) = -\mathcal{R}(A \wedge C)$.

\emph{(a).} The torsion equation at $a = \nabla\phi$ gives $D \wedge K = \hb^{T}(\nabla \wedge \nabla\phi) = 0$, so by (1) $A \cdot \mathsf{B}(C) - C \cdot \mathsf{B}(A) = (A \wedge C) \cdot (D \wedge K) = 0$. Next $K \cdot \mathsf{B}(A) = K \cdot (A \cdot D\,K) = \tfrac12\,\hb(A) \cdot \nabla\,Q(K) = 0$. By symmetry $A \cdot \mathsf{B}(K) = K \cdot \mathsf{B}(A) = 0$ for every $A$, so $\mathsf{B}(K) = 0$. Along $K$, $\hb(K) \cdot \nabla\phi = K \cdot \hb^{T}(\nabla\phi) = Q(K) = 0$: the curves of $K$ lie in the hypersurfaces of constant $\phi$.

\emph{(b).} $K \cdot D(\mathsf{B}(A)) = K \cdot D(A \cdot D\,K) = A \cdot D(K \cdot D\,K) + [K \cdot D, A \cdot D]K$. The first term vanishes by (a). By (3) the second is $-\mathcal{R}(K \wedge A) \times K + (K \cdot D\,A - A \cdot D\,K) \cdot D\,K = -\mathcal{R}(K \wedge A) \cdot K + \mathsf{B}(K \cdot D\,A) - \mathsf{B}(\mathsf{B}(A))$, the commutator product of a bivector with a vector being their inner product.

\emph{(c).} By (a), $\mathsf{B}$ annihilates $K$ and has image in the orthogonal complement of $K$. Take physical vectors $m_{1}, m_{2}$ orthogonal to $K$ and to each other with $Q(m_{i}) = -1$, and a vector $L$ orthogonal to them with $K \cdot L \neq 0$. In the basis $K, m_{1}, m_{2}, L$ no image has a component along $L$ (it would not be orthogonal to $K$) and the image of $K$ is zero, so $\mathsf{B}(m_{i}) = \mathsf{S}(m_{i}) + \alpha_{i}K$ with $\mathsf{S}$ a map of the span of $m_{1}, m_{2}$, and $\mathsf{B}^{2}(m_{i}) = \mathsf{B}(\mathsf{S}(m_{i})) = \mathsf{S}^{2}(m_{i}) + (\cdots)K$: the traces of $\mathsf{B}$ and $\mathsf{B}^{2}$ are those of $\mathsf{S}$ and $\mathsf{S}^{2}$. $\mathsf{S}$ is symmetric for the product on the screen, which is definite; so $\mathsf{S} = \tfrac12\theta + \sigma$ with $\sigma$ symmetric and trace-free, $\tr\mathsf{S}^{2} = \tfrac12\theta^{2} + \tr\sigma^{2}$, and $\tr\sigma^{2} = \ip{\sigma}{\sigma} \geq 0$. The left side of (b) is linear in the value of $A$ at the point (replacing $A$ by $fA$ multiplies it by $f$), so it defines a map $K \cdot D\,\mathsf{B}$; on a constant frame it is the derivative of the components of $\mathsf{B}$ along $\hb(K)$ plus the commutator of $\mathsf{B}$ with $X \mapsto \Omega(\hb(K)) \times X$, which has no trace. So its trace is $d\theta/d\lambda$. The trace of the curvature term: $\mathcal{R}(K) = \partial_{A} \cdot \mathcal{R}(A \wedge K)$, so $K \cdot \mathcal{R}(K) = K \cdot \bigl(\partial_{A} \cdot \mathcal{R}(A \wedge K)\bigr) = \partial_{A} \cdot \bigl(\mathcal{R}(A \wedge K) \cdot K\bigr) = -\partial_{A} \cdot \bigl(\mathcal{R}(K \wedge A) \cdot K\bigr)$, by the identity for $x \cdot (\mathcal{B} \cdot y)$ above. Taking the trace of (b): $d\theta/d\lambda = -\tr\mathsf{B}^{2} + K \cdot \mathcal{R}(K) = -\tfrac12\theta^{2} - \ip{\sigma}{\sigma} + K \cdot \mathcal{R}(K)$. For the shear, (b) restricts to the screen: for $A$ orthogonal to $K$, $K \cdot (K \cdot D\,A) = \hb(K) \cdot \nabla(K \cdot A) - (K \cdot D\,K) \cdot A = 0$, and $K \cdot D\,K = 0$. On the screen $\mathsf{S}^{2} = \tfrac14\theta^{2} + \theta\sigma + \sigma^{2}$ and, in two dimensions, a symmetric trace-free map has $\sigma^{2} = \tfrac12\ip{\sigma}{\sigma}$ times the identity; so the trace-free part of $\mathsf{S}^{2}$ is $\theta\sigma$, and the trace-free part of (b) is $K \cdot D\,\sigma = -\theta\sigma - \widehat{\mathcal{C}}$ with $\mathcal{C}(a, b) = b \cdot (\mathcal{R}(K \wedge a) \cdot K) = (K \wedge b) \cdot \mathcal{R}(K \wedge a)$. With the Weyl map of the Definition, $\mathcal{R}(K \wedge a) = \mathcal{W}(K \wedge a) + P(K) \wedge a + K \wedge P(a)$, $P$ the Schouten map, and for $a, b$ orthogonal to $K$: $(K \wedge b) \cdot (P(K) \wedge a) = (K \cdot a)(b \cdot P(K)) - (K \cdot P(K))(b \cdot a) = -(K \cdot P(K))\,a \cdot b$ and $(K \wedge b) \cdot (K \wedge P(a)) = (K \cdot P(a))(b \cdot K) - Q(K)(b \cdot P(a)) = 0$: the two forms differ by a multiple of the product. The check builds a field $(\hb, \Omega)$ that is not weak, with $\Omega$ solved exactly from the torsion equation, and verifies (a), the Riccati equation on the four frame vectors, the trace identity, the focusing equation with $d\theta/d\lambda$ computed independently, the decomposition of $\tr\mathsf{B}^{2}$, the trace-free part of $\mathsf{S}^{2}$ and the Weyl statement, all in exact rational arithmetic.

\emph{(iv).} The inverse and the adjoint of a symmetric map are symmetric, so $\hb^{T} = \hb$, and $\hb^{-1}(n) = n$ gives $\hb(n) = n$. Then $K = \hb^{T}(\nabla(n \cdot x)) = \hb^{T}(n) = n$, $Q(K) = 0$, and (iii) applies; $\hb(K) = n$, so $K \cdot D = D_{n}$, which is $\partial_{n}$ on scalars, and $\lambda$ is the parameter of (iii)(a). For constant background vectors $a, b$ put $g_{a} = \hb^{-1}(a)$, so $g(a, b) = g_{a} \cdot g_{b}$ and $\partial_{n}g(a, b) = (K \cdot D\,g_{a}) \cdot g_{b} + g_{a} \cdot (K \cdot D\,g_{b})$. Lemma (2) with $A = K$, $C = g_{a}$, whose background images $n$ and $a$ are constant, gives $K \cdot D\,g_{a} = g_{a} \cdot D\,K = \mathsf{B}(g_{a})$. With the symmetry of $\mathsf{B}$, $\partial_{n}g(a, b) = 2\,g_{a} \cdot \mathsf{B}(g_{b})$. At first order $g_{a} = a + s(a)$, the derived metric is $a \cdot b + 2s(a) \cdot b$, so $h = 2s$ and $h(n) = 0$; $\mathsf{B}$ is of first order, so $g_{a} \cdot \mathsf{B}(g_{b}) = a \cdot \mathsf{B}(b)$ at this order and $\mathsf{B} = \tfrac12\partial_{n}h$. The screen of $K = n$ is represented by the transverse plane, on which $\mathsf{S}$ is the transverse part of $\tfrac12\partial_{n}h$; in the gauge the trace of $h$ is the trace of its transverse part (the reciprocals of $n, \bar n$ are $\tfrac12\bar n, \tfrac12 n$, and $h(n) = 0$), so $\theta = \tfrac12\partial_{n}\tau$ and $\sigma = \tfrac12\partial_{n}h^{TT}$. At first order the focusing equation of (iii) is $\partial_{n}\theta = n \cdot \mathcal{R}(n)$, the quadratic terms dropping, and $n \cdot \mathcal{R}(n) = n \cdot \mathcal{E}[h](n)$ since $Q(n) = 0$: the identity $n \cdot \mathcal{E}[h](n) = \tfrac12\partial_{n}^{2}\tau$ of (ii), obtained there from the stored formula alone. For the curvature, with $h(n) = 0$ and transverse $u$: $(n \wedge u) \cdot p = n\,(u \cdot p) - u\,(n \cdot p)$, so $h((n \wedge u) \cdot p) = -(n \cdot p)h(u)$ and $\mathcal{R}_{p}[h](n \wedge u) = \tfrac12(n \cdot p)\,p \wedge h(u)$; then $n \cdot (p \wedge h(u)) = (n \cdot p)h(u) - (n \cdot h(u))p = (n \cdot p)h(u)$, giving $n \cdot \mathcal{R}[h](n \wedge u) = \tfrac12\,\partial_{n}^{2}h(u)$, and $(n \wedge b) \cdot \mathcal{R}[h](n \wedge a) = -b \cdot \bigl(n \cdot \mathcal{R}[h](n \wedge a)\bigr) = -\tfrac12\,\partial_{n}^{2}h(a) \cdot b$, with trace-free part $-\partial_{n}\sigma$ as a form: $\partial_{n}\sigma = -\widehat{\mathcal{C}}$, the shear equation of (iii) at first order. The curvature vanishes on displacements (Theorem \claimref{synthesis.weyl-flat-variations-carry-only-the-stress-trace-and-the-whole-stress-clause-is-the-trace-free-channel}(i)), hence the statement on the gauge. \emph{The stored formula.} Let $\hb^{-1} = \id + s$, $s = \tfrac14(p \cdot x)^{2}h_{0}$. For constant $a, b$, Lemma (2) reads $D_{a}g_{b} = D_{b}g_{a}$, that is $(a \cdot \nabla)s(b) + \Omega(a) \cdot (b + s(b)) = (b \cdot \nabla)s(a) + \Omega(b) \cdot (a + s(a))$. If $\Omega(a) \cdot b = \Omega(b) \cdot a$ for all $a, b$ then $\Omega = 0$ (the form $c \cdot (\Omega(a) \cdot b)$ is symmetric in $a, b$ and antisymmetric in $b, c$, hence zero), so the equation fixes $\Omega$, and to first order in $s$ its solution is $\Omega(a) = -\nabla \wedge s(a)$, the derivative acting on $s$: indeed $(\nabla \wedge s(a)) \cdot b = \dot\nabla(\dot s(a) \cdot b) - (b \cdot \nabla)s(a)$, and the gradient terms cancel between $a$ and $b$ by the symmetry of $s$. Near $x = 0$, $s$ is of second order in $x$ and $\Omega(a) = -\tfrac12(p \cdot x)\,p \wedge h_{0}(a)$ up to terms of third order, so at $x = 0$, where $\Omega$ vanishes and $\hb$ is the identity, $R(a \wedge b) = (a \cdot \nabla)\Omega(b) - (b \cdot \nabla)\Omega(a) = -\tfrac12\,p \wedge h_{0}\bigl((a \cdot p)b - (b \cdot p)a\bigr) = \tfrac12\,p \wedge h_{0}((a \wedge b) \cdot p)$, and $\mathcal{R}(a \wedge b) = -R(a \wedge b) = \mathcal{R}_{p}[h_{0}](a \wedge b)$. The check verifies the curvature identities on four generic maps of the gauge and their invariance under a displacement that takes the map out of the gauge, the relation $\tfrac12\partial_{n}g(a, b) = \hb^{-1}(a) \cdot \mathsf{B}(\hb^{-1}(b))$ on the field that is not weak (with (iii)), and the stored formula against the curvature computed from the gauge fields on all six planes.

\emph{(v).} A map $E_{p}$ commuting with the rotors is the same thing as a form $F(k, h, p) = \ip{k}{E_{p}[h]}$, linear in the symmetric maps $k$ and $h$ and quadratic in $p$, unchanged when one rotor acts on all three; the rotors being generated by the bivectors, this holds exactly when $F$ is annihilated by each unit bivector $G = e_{\mu}e_{\nu}$ acting on every vector by the commutator product, $v \mapsto G \times v$. The coefficients of $F$ on triples of symmetric pairs of frame vectors number $10^{3}$, and the six bivectors give a linear system in them with integer coefficients. Its rank modulo the prime $1000003$ is $995$ (computed); a minor that is nonzero modulo a prime is nonzero, so the rational rank is at least $995$ and the space of solutions has dimension at most five. The five displayed maps are built from $p$ and $h$ with the product of the algebra, so they commute with the rotors, and they are linearly independent (computed on four generic arguments). Hence the space has dimension exactly five. On a displacement $h = (a \mapsto (a \cdot p)\xi + (a \cdot \xi)p)$: $\tau = 2\,p \cdot \xi$, $h(p) = Q(p)\xi + (p \cdot \xi)p$, $p \cdot h(p) = 2Q(p)(p \cdot \xi)$, so the five maps take the values $Q(p)h$, $2Q(p)(p \cdot \xi)\id$, $2Q(p)(p \cdot \xi)\id$, $2(p \cdot \xi)\Pi_{p}$ and $Q(p)h + 2(p \cdot \xi)\Pi_{p}$, with $\Pi_{p}(a) = (a \cdot p)p$. A combination with coefficients $c_{1}, \dots, c_{5}$ takes the value $(c_{1} + c_{5})Q(p)h + 2(c_{2} + c_{3})Q(p)(p \cdot \xi)\id + 2(c_{4} + c_{5})(p \cdot \xi)\Pi_{p}$, and the three maps $h$, $\id$, $\Pi_{p}$ are independent for generic $p, \xi$; so it vanishes on displacements exactly when $c_{5} = -c_{1}$, $c_{3} = -c_{2}$, $c_{4} = c_{1}$, that is, when it is $2c_{1}\mathrm{Ric}_{p} + c_{2}\,\id\,\tr\mathrm{Ric}_{p}$, by the formula for the Ricci map above. Of the pairings $\ip{h_{1}}{\cdot\,[h_{2}]}$ of the five maps, those of the first, second and fifth are symmetric in $h_{1}, h_{2}$ ($Q(p)\ip{h_{1}}{h_{2}}$, $Q(p)\tau_{1}\tau_{2}$, $2\,h_{1}(p) \cdot h_{2}(p)$), while the third gives $\tau_{1}\,p \cdot h_{2}(p)$ and the fourth $\tau_{2}\,p \cdot h_{1}(p)$; so the antisymmetric part of the two-parameter family is $(c_{1} + c_{2})\bigl(\tau_{2}\,p \cdot h_{1}(p) - \tau_{1}\,p \cdot h_{2}(p)\bigr)$, which vanishes exactly when $c_{2} = -c_{1}$: the family $2c_{1}\mathcal{G}_{p}$. Likewise $\mathrm{Ric}_{p}[h](p) = \tfrac12(\tr\mathrm{Ric}_{p}[h])\,p$, so the value at $p$ of a member of the family is $(c_{1} + c_{2})(\tr\mathrm{Ric}_{p}[h])\,p$, zero for all $h$ under the same condition. The last sentence of the clause is the passage from symbols to operators: an equation of the kind (L) has a symbol of the kind counted, (D) is the vanishing of the symbol on displacements, and a constant-coefficient operator is the stationarity condition of a quadratic functional exactly when its symbol is symmetric for the pairing. The check carries out the count and verifies each step on generic arguments: the rank $3$ of the conditions from displacements, the rank $4$ with symmetry and with conservation, and the identity of the surviving line with the stored Einstein map.

\emph{(vi).} Along $\hb = \id + t\,\delta$ the derived metric $\hb^{-1}(a) \cdot \hb^{-1}(b)$ is $a \cdot b - 2t\,\delta(a) \cdot b$ to first order for symmetric $\delta$, so $h = -2t\,\delta$, and the derivative of the matter density, $\ip{\delta}{T}$ by clause (v) of the cited theorem, is $-\tfrac12\ip{h}{T}$ per unit of $h$. By (v), under (L), (D), (V) the functional of the field is $c\int\ip{h}{\mathcal{E}[h]}$, and with the response it is stationary exactly when $2c\,\mathcal{E}[h] = \tfrac12 T$, $\mathcal{E}$ being symmetric for the pairing. On a wave of momentum $k$ the operator $\mathcal{E}$ acts as $-\mathcal{G}_{k}$ (two derivatives, each the momentum times a unit of square $-1$), so $\mathcal{G}_{k}[h] = -T/4c$. For a source at $k$, $T(k) = 0$, the Ricci formula at $h = T - \tfrac12\,\id\,\tr T$ (trace $-\tr T$, value $-\tfrac12(\tr T)k$ at $k$) gives $\mathrm{Ric}_{k} = \tfrac12 Q(k)(T - \tfrac12\,\id\,\tr T)$ with trace $-\tfrac12 Q(k)\tr T$, so $\mathcal{G}_{k}[T - \tfrac12\,\id\,\tr T] = \tfrac12 Q(k)\,T$, and the field of the source is $h = -\tfrac{1}{2c}(T - \tfrac12\,\id\,\tr T)/Q(k)$ up to a displacement. On the solution the cross term of the functional and the responses of two sources is $-\tfrac12\ip{h[T_{2}]}{T_{1}}$, so their interaction is $\tfrac12\ip{T_{1}}{h[T_{2}]} = -\tfrac{1}{4c}B(T_{1}, T_{2})/Q(k)$, displacements pairing to zero with a source. For two bodies at rest, $T_{i}(a) = m_{i}(a \cdot e_{0})e_{0}$: $B = m_{1}m_{2} - \tfrac12 m_{1}m_{2} = \tfrac12 m_{1}m_{2}$, and a static momentum has $Q(k) = -|\mathbf{k}|^{2}$, so the interaction is $m_{1}m_{2}/(8c\,|\mathbf{k}|^{2})$. By (N) this is the coupling $-Gm_{1}m_{2}/r$ of Proposition \claimref{field-theory.operator-newtonian-coupling}, whose transform is $-4\pi G m_{1}m_{2}/|\mathbf{k}|^{2}$ (the transform of $1/r$ is $4\pi/|\mathbf{k}|^{2}$, standard); equality holds exactly for $c = -1/32\pi G$. The Proposition takes that coupling as given, so (N) is the identification that gives $G$ its meaning, not a derivation of Newton's law. Then $2c\,\mathcal{E}[h] = \tfrac12 T$ is $\mathcal{E}[h] + 8\pi G\,T = 0$, the field of a source is $16\pi G\,(T - \tfrac12\,\id\,\tr T)/Q$ and the interaction $8\pi G\,B/Q$. A functional $c\int\ip{h}{\mathcal{E}[h]}$ has, by (i), the symplectic current $2|c|\,\omega$ up to orientation, here $\omega/16\pi G$. On maps transverse and trace-free at their momentum the identity above gives $\mathcal{G}_{p} = \tfrac12 Q(p)$, so $\mathcal{E} = \tfrac12\nabla^{2}$ there and the functional has density $\tfrac{1}{64\pi G}\ip{\partial h}{\partial h}$ up to a divergence, of the sign that makes the energy positive (the pairing is positive on these maps). With a source, each component then obeys $Q(p)\,h_{i} = 8\pi G\,\ip{\epsilon_{i}}{T}$, and the exchange $\tfrac12\sum_{i}\ip{T_{1}}{\epsilon_{i}}h_{i}[T_{2}]$ has residue $8\pi G \cdot \tfrac12\sum_{i}\ip{T_{1}}{\epsilon_{i}}\ip{T_{2}}{\epsilon_{i}}$ at the cone, the form of clause (iv) of Theorem \claimref{synthesis.first-order-gravitational-field-has-two-free-directions-at-null-momentum-and-four-constraints}; the check confirms on nine pairs of conserved sources at the null momentum $n$ that this is $8\pi G\,B$ there and that no other constant in the density gives it. So the Newtonian coupling, the stored residue and the sign of the energy agree on one constant. The check solves the equation $\mathcal{G}_{k}[h] = -T/4c$ for the field of a body at rest by exact elimination, forms the interaction from the value of the functional on the solutions, and finds $-4\pi G m_{1}m_{2}/|\mathbf{k}|^{2}$ exactly at $c = -1/32\pi G$ and not at neighbouring values; it does the same at a momentum off the cone for sources with transverse trace-free parts. On $N$, by (ii), the flux of $\omega/16\pi G$ is $\tfrac{1}{16\pi G}\int\tfrac12(\ip{\partial_{n}h^{TT}}{h'^{TT}} - \ip{h^{TT}}{\partial_{n}h'^{TT}})$, and with $x = \lambda n + y$ the measure on $N$ is $d\lambda\,d^{2}y$ (in the labels $n \cdot x$, $\bar n \cdot x = 2\lambda$, $y$ the volume is $\tfrac12\,d(n \cdot x)\,d(\bar n \cdot x)\,d^{2}y$). Integrating by parts along the generators and choosing the orientation, $\varpi(h, h') = \tfrac{1}{16\pi G}\int\ip{h^{TT}}{\partial_{n}h'^{TT}} = \tfrac{1}{8\pi G}\sum_{i}\int h_{i}\,\partial_{n}h_{i}'$. For $P = \tfrac{1}{32\pi G}\int\ip{\partial_{n}h^{TT}}{\partial_{n}h^{TT}}$, $dP(h') = \tfrac{1}{16\pi G}\int\ip{\partial_{n}h^{TT}}{\partial_{n}h'^{TT}} = \varpi(\partial_{n}h, h')$, so $P$ generates $\lambda \mapsto \lambda + t$, it is nonnegative, and with $\sigma = \tfrac12\partial_{n}h^{TT}$ it is $\tfrac{1}{8\pi G}\int\ip{\sigma}{\sigma}$. The opposite orientation of $\varpi$ would make the generator of the future translations $-P$; the orientation is fixed by positive null energy, as for the photon.

\emph{(vii).} By orthogonality of $\cos k\lambda$ and $\sin k\lambda$ on $[0, 2\pi]$, each square integrating to $\pi$, $\int_{0}^{2\pi}h_{i}\,\partial_{\lambda}h_{i}'\,d\lambda = \pi\sum_{k}k(a_{ik}b_{ik}' - b_{ik}a_{ik}')$ and $\int_{0}^{2\pi}(\partial_{\lambda}h_{i})^{2}d\lambda = \pi\sum_{k}k^{2}(a_{ik}^{2} + b_{ik}^{2})$; with the cell area $s$ this gives the stated $\varpi$ and $P$, the factors $\pi$ cancelling against those of $8\pi G$ and $16\pi G$. The translation $\lambda \mapsto \lambda + t$ moves $(a_{ik}, b_{ik})$ with velocity $(kb_{ik}, -ka_{ik})$, and $\varpi(X, Y) = dP(Y)$ for every $Y$ exactly for this $X$: $\tfrac{s}{8G}k(X_{a}Y_{b} - X_{b}Y_{a}) = \tfrac{s}{8G}k^{2}(aY_{a} + bY_{b})$. The corresponding bracket is $\{a_{ik}, b_{jl}\} = 8G\,\delta_{ij}\delta_{kl}/(sk)$, with $\{a_{ik}, P\} = kb_{ik}$. Under $[\,,\,] = U\{\,,\,\}$ with a central unit $U$, $U^{2} = -1$, $U^{*} = -U$, and $a$, $b$ self-adjoint: $[a_{ik} + Ub_{ik}, a_{ik} - Ub_{ik}] = -2U[a_{ik}, b_{ik}] = 16G/(sk)$, so $[\alpha_{ik}, \alpha_{ik}^{*}] = (s/16G)k^{2} \cdot 16G/(sk) = k$; modes with different $(i, k)$ commute, and $[a + Ub, a + Ub] = 0$. From $P = \tfrac{s}{16G}\sum k^{2}(a^{2} + b^{2})$, $[P, a_{ik}] = -Ukb_{ik}$ and $[P, b_{ik}] = Uka_{ik}$, so $[P, a_{ik} + Ub_{ik}] = -Ukb_{ik} + U^{2}ka_{ik} = -k(a_{ik} + Ub_{ik})$. Replacing $U$ by $-U$ in both the rule and the modes gives $+k$ again. The unit is $I_{4}$ by the unit principle for the scalars of an abstract mode space (the conventions, clause (U1)), which is the unit of Premise (TT). Then $k^{-1/2}\alpha_{ik}$ and its adjoint satisfy the canonical relation and are represented as in Proposition \claimref{field-theory.fock-space-maxwell-modes}(1). In a state of $N$ quanta, $\alpha^{*}\alpha + \alpha\alpha^{*} = (s/8G)k^{2}(a^{2} + b^{2})$ has the value $(2N + 1)k$, so $a^{2} + b^{2} = 8G(2N + 1)/(sk)$, whatever the unit of length, $s$ and $G$ being measured in the same one; in units $\hbar = c = 1$, $G$ is the square of the Planck length. The closed generator is the periodic setting of discretised light-cone quantisation, a compactification of a null direction of the background, used here for the mode count only. The check verifies the two integrals on generic coefficients at three levels, the translation, and the commutators with the unit held as a number of square $-1$, for both orientations.

\emph{(viii).} $R = \cos(\phi/2) + e_{1}e_{2}\sin(\phi/2)$ commutes with $e_{0}$ and $e_{3}$, so it fixes $n$ and $\bar n$; $e_{1}$ and $e_{2}$ anticommute with $e_{1}e_{2}$, so $Re_{1}\tilde R = e_{1}\tilde R^{2} = \cos\phi\,e_{1} + \sin\phi\,e_{2}$ and $Re_{2}\tilde R = -\sin\phi\,e_{1} + \cos\phi\,e_{2}$. The rotated field $h'(a) = Rh(\tilde RaR)\tilde R$ of $h = h_{+}\epsilon_{+} + h_{\times}\epsilon_{\times}$ is $h_{+}\epsilon_{+}' + h_{\times}\epsilon_{\times}'$ with $\epsilon'$ built on the rotated vectors, and $\epsilon_{+}' = \cos 2\phi\,\epsilon_{+} + \sin 2\phi\,\epsilon_{\times}$, $\epsilon_{\times}' = -\sin 2\phi\,\epsilon_{+} + \cos 2\phi\,\epsilon_{\times}$ (the doubling of the cited theorem, proved under its clause (i)), so $h_{+}' = \cos 2\phi\,h_{+} - \sin 2\phi\,h_{\times}$, $h_{\times}' = \sin 2\phi\,h_{+} + \cos 2\phi\,h_{\times}$ and $h_{+}' \pm Uh_{\times}' = (\cos 2\phi \pm U\sin 2\phi)(h_{+} \pm Uh_{\times})$; the trace is invariant. The modes are linear in the components with the unit $I_{4}$, so $\beta_{\pm,k}$ are multiplied by $e^{\pm 2I_{4}\phi}$, and by (vii) both by $e^{-I_{4}kt}$ under the null translation. With $\beta_{-,l}^{*} = 2^{-1/2}(\alpha_{+l}^{*} + I_{4}\alpha_{\times l}^{*})$ and $[\alpha_{ik}, \alpha_{jl}^{*}] = k\delta_{ij}\delta_{kl}$: $[\beta_{+,k}, \beta_{-,l}^{*}] = \tfrac12(k + I_{4}^{2}k)\delta_{kl} = 0$ and $[\beta_{\pm,k}, \beta_{\pm,l}^{*}] = \tfrac12(k - I_{4}^{2}k)\delta_{kl} = k\,\delta_{kl}$. The check uses the unnormalised rotor $2 + e_{1}e_{2}$, for which $\cos\phi = \tfrac35$ and $\cos 2\phi = -\tfrac{7}{25}$.

\emph{Reading.} For null $K$ the scalar part of the Einstein map drops, $K \cdot \mathcal{G}(K) = K \cdot \mathcal{R}(K)$, so a field equation $\mathcal{G}(a) + 8\pi G\,T(a) = 0$ gives $K \cdot \mathcal{R}(K) = -8\pi G\,K \cdot T(K)$; substituted in the focusing equation of (iii) this is the displayed form, in which $\ip{\sigma}{\sigma}$ and $8\pi G\,K \cdot T(K)$ stand with the same sign, and by (vi) the first is $8\pi G$ times the null energy density of the two components. The shear being of first order in $h$, the comparison is at second order in $h$, the leading one.
